%!TEX root = ../../main.tex
\section{문제}
\label{sec:intro-problem}

교전, 곧 한타와 소규모 교전을 아우른 단위의 결과를 어떤 양으로 정의할 것인가. 본 논문은 이 물음에 두 모형으로 답하고, 둘째 모형은 첫째 모형 위에 설계한다. 첫째 모형은 경기의 승패를 예측한다. 어느 시점의 상태를 넣으면 그 경기가 어느 쪽의 승리로 끝날 확률을 낸다. 교전의 결과는 이 모형이 매긴 추정 승률이 그 구간에서 오르는지이며, 오르지 않은 경우, 곧 내리거나 변하지 않은 경우는 같은 쪽에 둔다. 본 논문은 이 증가 여부를 전략적 가치 개선(strategic value improvement, SVI)이라 부른다. 둘째 모형은 그 증가 여부를 예측한다. 이 정의는 킬 교환만으로 결과를 정하는 대신, 경기 승패와 관련된 상태 변화를 하나의 기준으로 요약한다. 그 구분이 지도 통제권을 얼마나 반영하는지는 첫째 모형이 상태를 어떻게 읽는가에 달려 있고, 구간 전후의 차이가 그 교전만의 영향이라고 보장하지는 않는다.

본 논문은 킬 기록으로 사후에 가려낸 교전을 대상으로, 첫 킬보다 15초 앞까지 공개된 상태만을 쓴다. 이 시각은 피해가 오가기 시작하기 전이라는 뜻이 아니라, 예측에 쓸 시점을 그렇게 정해 둔 것이다. 예측하는 것은 고정한 첫째 모형이 평가한 승률의 증가 여부이고, 교전이 일어날지 또는 그 교전이 승리를 일으켰는지는 아니다. 그러므로 이 문제는 시작 승률과 경기 시각을 넘어서, 선택한 사전 상태가 얼마나 더 예측하는지를 재는 것이다. 예측 시점과 결과 시점의 정의는 제 \ref{ch:engagement} 장에, 승률 추정치와 증가 여부의 정의는 제 \ref{ch:value} 장에 있으며, 절차와 결과의 요약은 제 \ref{sec:intro-approach} 절에 둔다.
